%% Copyright 2026
%% Licensed under Creative Commons Attribution–ShareAlike 4.0 International (CC BY-SA 4)

\documentclass[12pt]{article}
\usepackage{titling}
\usepackage{listings}
\usepackage{xcolor}
%% \usepackage{fontspec}

\usepackage{sourcecodepro}

\newcommand{\NULL}{{\texttt{NULL}}}%
\newcommand{\cvolatile}{{\texttt{volatile}}}%
\newcommand{\cdefine}{{\texttt{\#define}}}%
\newcommand{\cstatic}{{\texttt{static}}}%
\newcommand{\cconst}{{\texttt{const}}}%
\newcommand{\cCC}{{\texttt{C/C++}}}%
\newcommand{\cC}{{\texttt{C}}}%

\lstset{ %
   numbers=left,   
   stepnumber=2,
   numbersep=5pt,
   numberstyle=\footnotesize,
   language=C++,
   basicstyle=\footnotesize\ttfamily,
   breaklines=true,
   frame=shadowbox,
}

\author{spinoff@protonmail.com}
\date{August 27, 2021}

\title{So, you want to be an \\embedded software engineer?}

\begin{document}
\begin{titlingpage}
\maketitle
\begin{abstract}
A lot of ink has been used to describe the requirements for
a good embedded software engineer.  This document focuses on the core
competency of an embedded software engineer. Good luck.
\end{abstract}
\par
\vskip1cm
{\textsc{version 1.0}}
\vskip5mm

\end{titlingpage}

\section{Introduction}

Where noted ($\ddagger$), these questions were stolen from an article
in Embedded Systems Journal in 2002.

When I first read the questions back then it occurred to me that some work needed to be
done within the company I was with at the time to enhance our test for prospective 
engineers.  I've dragged this list around with me since.

The nature of these questions is not to be misinterpreted -- a candidate should know
how to solve all these problems. If they cannot, they just cannot be hired.

I could see perhaps a new graduate not having enough experience but there's a certain
cut-off where the answers should just flow out rapidly.  I'll mark in the document
where that cut-off should be.


Here we go! 

\section{The C-Test}

\begin{enumerate}

\item $\ddagger$ Using the \cdefine\ directive, how would you declare a manifest
constant that is the number of seconds in a year (Disregard leap year in your
answer).
\par
What we're looking for here is something like this. Your candidate answer may
vary but the point is clear:
\par
\begin{lstlisting}
#define DAYS_PER_YEAR (365)
#define HOURS_PER_DAY (24)
#define MINUTES_PER_HOUR (60)
#define SECONDS_PER_MINUTE (60)
#define SECONDS_PER_DAY (SECONDS_PER_MINUTE * MINUTES_PER_HOUR * HOURS_PER_DAY)
#define SECONDS_PER_YEAR (SECONDS_PER_DAY * DAYS_PER_YEAR)
\end{lstlisting}

\begin{itemize}
\item Did they use some logical build-up of the answer? or did they just write:
\begin{lstlisting}
#define SECONDS_PER_YEAR (31536000)
\end{lstlisting}
Why is that a bad idea?  What's good about it? 
\item Did they name the manifest constant clearly?
\item Did they enclose the values in {\texttt{( )}}?
\item Did they mark the values with {\texttt{U}} (for unsigned) ?  Does it matter? Why or why not?
\item Why do we use \cdefine\  rather than:
\begin{lstlisting}
const unsigned int SECONDS_PER_YEAR = 315360000U;
\end{lstlisting}
What's the impact? What's the difference between that and using \cdefine?
\end{itemize}

\item $\ddagger$ Write the ``standard'' {\texttt{MIN}} macro such that the macro takes two arguments
and resolves to the value of the minimum of the arguments.
\par
We're looking for something like this:
\par
\begin{lstlisting}
#define MIN(x,y) (((x) < (y)) ? (x) : (y))
\end{lstlisting}
\par
Did they guard the values with parentheses?

\par
So, they probably did OK and wrote a macro like that.. So let's dig deeper:

Comment on this code:
\begin{lstlisting}
int a;
int b;
int *pa = &a;
// ...
int r = MIN(*pa++, b);
\end{lstlisting}

There's a lot to unpack here.  Here's what the compiler sees:
\begin{lstlisting}
int r = (((*pa++) < (b)) ? (*pa++) : (b));
\end{lstlisting}

What is wrong with this scenario? Did they find the possibility that whatever {\texttt{pa}} is pointing to is incrementing
twice?  The side effects are bad.  Did they comment/find this problem?

\item $\ddagger$ What is the purpose of the pre-processor directive {\texttt{\#error}} ?

They either know this or they do not.  If they don't, then they have likely not done any embedded software
development that had significant use of headers which differentiate a platform, or a compiler, or something
else that would be correct or incorrect code for a scenario.

If that's true -- they've never seen {\texttt{\#error}}, I'm not sure I want them working on embedded software.
It's not that the topic is difficult to understand. It's indication that their experience is too light. And,
we're not looking for people barely qualified for embedded software.

\item $\ddagger$ Infinite loops often arise in embedded systems. How do you code an infinite loop?

There are just a few ways to do this, and they are all correct.  But which one they pick tells
a lot about background.

Most go this way:

\begin{lstlisting}
while(1)
{
   do_whatever();
}
\end{lstlisting}

A few go this way:
\begin{lstlisting}
for(;;) 
{
   do_whatever();
}
\end{lstlisting}

And even smaller subset might go this way:
\begin{lstlisting}
foo:
   do_whatever();
   goto foo;
\end{lstlisting}


Now, the first answer usually indicates a very correct answer that works everywhere that {\texttt{1}} means true.
In some languages like {\texttt{C\#}}, this would not be the case.  In \cCC\ it will work fine.
This candidate would fit in fine and would be able to survive anywhere.

\par
The next one -- the {\texttt{for}} loop is also OK.  It's a little pedantic.  People will understand it.  But it's
not what we hope to see.  It works, it's not incorrect.  But the syntax isn't clean in meaning by reading it.
\par

The last one -- {\texttt{goto}} ... Now this is a two edged sword.  It would be odd to get this answer
{\textit{unless}} they were an assembly language developer at heart.  That's good and bad.  It's good that
they are experienced at low level software.  It's bad because we haven't discovered if they know enough about \cC.  If you get this answer -- be careful with this candidate to fully understand 
their \cC\ experience.


\item $\ddagger$ Using the symbol {\texttt{x}}, give the declarations for the following:
\begin{itemize}
\item An integer.
\item A pointer to an integer.
\item A pointer to a pointer to an integer.
\item An array of ten integers.
\item An array of ten pointers to integers.
\item A pointer to an array of ten integers
\item A pointer to a function that takes an integer argument  and returns an integer.
\item An array of ten pointers to functions that each take an integer argument and each return an integer.
\end{itemize}

What we're looking for is quite simple.  And, frankly if this question is missed, the interview is over.

Here's what we want to see:
\begin{itemize}
\item An integer. {\texttt{int x;}}
\item A pointer to an integer. {\texttt{int* x;}}
\item A pointer to a pointer to an integer. {\texttt{int** x;}}
\item An array of ten integers. {\texttt{int x[10];}}
\item An array of ten pointers to integers. {\texttt{int* x[10];}}
\item A pointer to an array of ten integers  {\texttt{int (*x)[10];}}  (this one may catch your candidate.)
\item A pointer to a function that takes an integer argument  and returns an integer.
  {\texttt{ int (*x)(int); }}
\item An array of ten pointers to functions that each take an integer argument and each
return an integer.  {\texttt{ int (*x[10])(int); }}
\end{itemize}

These are the kinds of questions that may seem nit-picky.  They are.  And some may argue that
it would require the candidate to a bit of preparation ahead of time to be sure they
grok all the ways to declare things.  If they weren't prepared for the interview to know this,
then what are they prepared for?

\item $\ddagger$ What are the uses of the keyword {\texttt{static}} ?

A surprising number of candidates will not get all three answers, and a few will
outright not understand the differences between the usages.

\begin{itemize}
\item Function scope:  ``A \cstatic\ automatic variable in a function
retains it's last value between function calls.''
\item Module scope:  ``A \cstatic\ variable in a file (aka the module)
is accessibly only from functions within the same module.  A localized global variable.''  The 
variable is private to the module.
\item Module scope:  ``A \cstatic\ function in a file (aka a module)
is accessible only from functions within the same module. A localized function.''  The function
is private to the module.
\end{itemize}

This is one of those questions that either ends the interview or lets it proceed.

Any embedded software engineer applying should know this answer.

\item $\ddagger$ What does the keyword \cconst\  mean?

The question was about the meaning, not about what it does.  Let's focus on
the meaning.  \cconst\ means {\textsc{read only}}.  It does not
nor has it ever meant ``constant''.  If the candidate says ``constant'',
the interview is over.

The details:  The meaning is that a \cconst\ thing cannot be changed
after it has been assigned.  It's (wait for it) {\textsc{read only}}.

\par
Comment on this code:
\begin{lstlisting}
const int x;
const int y = 10;
// ...
x = 11;
\end{lstlisting}

\par
Here we're being crystal clear:  We are assigning {\texttt{y = 10}} and after that, the
value of {\texttt{y}} cannot be changed. It's read-only.  Immutable.
It does not mean constant. Here's why:

\begin{lstlisting}
const unsigned int* preg = (unsigned int*)(REGISTER_BASE_ADDRESS);
// ...

unsigned int tmp = *preg;
while ( tmp & 0x10 )
{
   // ...
   tmp = *preg;
}
\end{lstlisting}

{\texttt{preg}} points to a {\texttt{const unsigned int}}.  The value can change
because it's pointing to (perhaps) some register.  We're not allowed to change
the value in the code, but surely the value can be changed by the hardware
that it represents.

It's read-only.  It's surely not constant.  We covered what a constant
was in the first question.

That's the difference.

This leads to a set of questions here that should be addressed.
Comment on what each of these declarations mean:

\begin{itemize}
\item  {\texttt{const int x;}}
\item {\texttt{int const x;}}
\item {\texttt{int* const x;}}
\item {\texttt{const int* x;}}
\item {\texttt{const int* const x;}}
\end{itemize}

The first two are the same.\\
The third is a \cconst\ pointer to a non-\cconst\ integer.
The pointer is read-only, the integer it points to is not read-only.\\
The fourth is a pointer to a \cconst\ integer.  The pointer
is not read-only.  The integer is read-only.\\
The fifth is a {\texttt{const}} (read-only) pointer to a {\texttt{const}}
(read-only) integer.  The pointer nor the integer it points to can be
modified by the code.

We hope the candidate brings up a final point -- the {\texttt{const}}-ness of
something helps the programmer and the user of the software interface understand
what the intent of the code is.  It helps to prevent actions to be compiled into
code that would cause unintended consequences or reveal a capability that doesn't
exist.  It's a hint to the programmer.  It's also a hint to the compiler and
linker for where items that are {\texttt{const}} are held in the final 
software product.  This differentiation is sometimes important.  Sometimes
it's very important.

An embedded software engineer knows this.  An API programmer doesn't care.
\end{enumerate}

\section{Getting Warmed Up}

By now you should have a good idea of the candidate's {\texttt{C}} capabilities.
But we still need to dig deeper.
\par
\begin{enumerate}
\item $\ddagger$ What does the keyword {\texttt{volatile}} mean?  Give three examples
of its use.

Again, let's focus on the meaning.  It technically means ``A value
that can change unexpectedly.''

That is the text-book definition in the context of {\texttt{C/C++}}.

It's not really appropriate to get into the {\textit{implementation details}}
of how a compiler will interpret the {\texttt{volatile}} keyword.  What is
important is the intent of the keyword and where/why it is used.

Have the candidate comment on this code:

\begin{lstlisting}
int square(volatile int* px)
{
   return *px * *px;
}
\end{lstlisting}

Follow up questions:
\begin{itemize}
\item Is there a defect with that code?  The answer is yes.
The argument is a {\texttt{volatile int}}. Specifically, it's
a pointer to a {\texttt{volatile int}}.  What the code does,
as written, is dereference the value of this {\texttt{volatile int}}
and multiply it by the value of this {\texttt{volatile int}}. 
The problem is with the value as the expression is evaluated.
It is not guaranteed that the value of the first {\texttt{*px}}
is the same as the value of the second {\texttt{*px}}.  
Whatever {\texttt{px}} pointed to can change as the expression
is evaluated.  Therefore it won't be guaranteed to compute the
square of the value {\texttt{*px}}.
\item Is there a fix to this?  The answer is yes.
\begin{lstlisting}
int square(volatile int* px)
{
   int x = *px;
   return x * x;
}
\end{lstlisting}
Capture the value at the onset of the function. Square it
(the captured value is not volatile), and then return the
result.
\end{itemize}

\item Can a ``variable'' in C be both {\texttt{volatile}}
and {\texttt{const}}?
Ask the candidate to comment on this code:
\begin{lstlisting}
const volatile int* preg;
\end{lstlisting}

The answer is that it is perfectly fine to use code like
this for the situation where the value that {\texttt{preg}}
points to (a {\texttt{const volatile}} will change
unexpectedly {\textit{and}} be read-only. A perfect situation
for a register value.

\par
So in a nutshell, let us not get wrapped around the axle on what the compiler
does when it encounters the \cvolatile\ keyword.  The take-home message is that
the compiler is being instructed to treat that value as something which will
change unexpectedly.  Even if it doesn't.  The compiler will treat it as it does, and
that means fetching the value from whatever source is done whenever the value is needed.
Yes, the literature will say ``it means that the compiler will not optimize away the
read/write to that variable.''  Yes, but the meaning of what the keyword is is more important.
What a compiler does (a vendor specific issue) is only just partially related to the answer.
\par
So in the example of {\texttt{square()}}, that put a light on the defect where we did not
capture the value once, and use the safe captured value to compute the square.

\item $\ddagger$ Embedded systems often require the manipulation of bits in
registers or other parameters.  Given a variable 
{\texttt{unsigned int x}}, show how to:
  \begin{itemize}
    \item Set bit 3 of {\texttt{x}} (the fourth bit)
    \item Clear bit 3 of {\texttt{x}}
  \end{itemize}

The candidate either knows this or they do not.  There is no
{\textit{half-way}} answer.

What we are looking for is this:

\begin{lstlisting}
#define SET_BIT(a,n) (a) |= (1 << (n))
#define CLEAR_BIT(a, n) (a) &= ~(1 << (n))
// ...
unsigned int y;
// ...
SET_BIT(y, 3);
CLEAR_BIT(y, 3);
\end{lstlisting}

If the candidate does not choose to write a general use macro, then
an answer like this is sufficient:
\begin{lstlisting}
y |= (1 << 3);
y &= ~(1 << 3);
\end{lstlisting}

Here we're just looking for is the operators {\texttt{|=, \&=, and \~}}.

A candidate who cannot get this answered isn't hired.


\item $\ddagger$ Comment on this code and explain what is the output (and why)?

\begin{lstlisting}
void foo(void)
{
    unsigned int a = 6;
    int b = -20;
    
    ( a + b  > 6 ) ? puts("> 6\n") :
                     puts("<= 6\n");
}
\end{lstlisting}

In just a few words, what we're looking for here is ``unsigned promotion''.
Right?  What happens when the arithmetic has a mixture of signed and
unsigned values?  The signed value is interpreted as an unsigned value.
The unsigned value from {\texttt{int b = -20}} is a very large
unsigned integer.  So the answer is {\texttt{> 6}}.

\end{enumerate}

\section{Downhill now}

At this point we are nearing the end of the test.  If the candidate
is demoralized then that's a bad sign. They should have been able
to whip through these questions so far without a sweat.

But now we're going to get more elaborate in our questions and stray
beyond the pure syntax and use of {\texttt{C}} in the embedded system
scenario.

\begin{enumerate}
\item $\ddagger$ {\texttt{typedef}} is frequently used in C to declare a complex
data structure.   It's also possible to use the pre-processor to do
something similar.  For instance, consider the following code:

\begin{lstlisting}
#define dPS struct s *

// vs.

typedef struct s* tPS;
\end{lstlisting}

The intent in both cases is to define {\texttt{dPS}} and
{\texttt{tPS}} to be pointers to a structure {\texttt{s}}.

Which method is preferred and why?

The answer is we prefer using {\texttt{typedef}}.  The compiler
will help us track down mis-use of the data type.  It also
allows us to define formal parameters to functions that
accept or return values of that given type.  Although not always
a motivation in embedded software, attempts to inject a bit more of 
abstraction at least in terms of definitions and usages helps the
reader understand what's being done in the code.  And further,
an abundance of macros used to define complex types will be difficult
to maintain for the next engineer on your project.  

Humans can read the code and made the associations in their mind
of what a syntax means.  A pre-processor definition to define a type
doesn't read-off the code easily.  Frankly, it's archaic.

A lot of the embedded software engineer candidates we come across
did not start and work in embedded software their entire career.
They had careers in other software (API-level) software that forced
them to be very abstract and elaborative of their code.

{\texttt{typedef}} is the right answer for all these reasons.

\item Although it is not common in an embedded system, there are RTOS
in the industry that do synthesize some dynamic memory allocation
routines.  So, explain the the problem of using dynamic memory
in an embedded software application.  Describe the issues. Start 
with this simple question -- explain the difference between heap, 
free-store, and the stack?

A follow up question is:   Can you explain why I would want to avoid
too many Tasks in an application that uses an RTOS?  How many is too many?
What else should I be concerned about when I have a lot of tasks in an
embedded software application?

\par

This is probably the most soft-ball question given and there aren't 
stark answers.  It's going to be a discussion.   And through the discussion
you (the interviewer) will pick up on things that they think and it will
help shape your perception of what the candidate really knows about
overall embedded software.

Let's take the question apart and see what happens:

\begin{itemize}
\item RTOS.  Can you get them to talk about what RTOS they have used?
The author likes to take a diversion here and level set what the 
candidate thinks an {\textit{Operating System}} is?  Here's what
should be in their answer:  An operating system provides capabilities to\ldots
   \begin{enumerate}
      \item Create and manage processes (Tasks, processes, etc\ldots)
      \item Manage memory (including handling allocation of and protection of memory)
      \item Handle I/O (otherwise known as device drivers to specific hardware)
      \item Filesystem (aka, persistent storage)
   \end{enumerate}
That's what you'd get from the first chapter of Tennenbaum's
{\underline{Operating Systems}} textbook.
But about RTOS -- what do they use in the RTOS?  Where do they put the
{\textsc{mutex}} and the {\textsc{semaphore}}?   
\item Memory.  Where does the code go?  What is the method that they use to
define where instructions and data are placed in the device?  (We're looking for 
{\textit{link files}}, {\textit{scatter files}}, etc\ldots).  But what
do those files do?  
\item Heap, stack and free-store.   A thing is allocated depending on how it's
instantiated.   Getting the candidate to talk about the difference and properties of
the heap, stack and free-store is an interesting test.   First, the chances of 
using any dynamic memory are suspect in an embedded system unless a type of RTOS
is even used.   We may not ever use one.  It may be bare-metal C, period.  But
if we are using an RTOS and we are dynamically allocating {\textit{anything}}
this question is important.

\begin{itemize}
   \item Stack.   The stack is where automatic variables are stored.
\begin{lstlisting}
void foo()
{
   int bar_on_stack;
   // etc...
   do_something();
   // etc...
}
\end{lstlisting} 
Access is fast since  the access is just an offset from the frame pointer.  The
space provided for the stack is transient and subject to the call to the function. The
compiler injects the instructions which sets the stack and frame pointer as the function 
call is made (those instructions to perform the linkage on entry and exit of the function).

For example, MC68000 assembly uses this:

\begin{lstlisting}
SUBR	 MOVEM.L  DO-D2,-(A7)
         ## code here
         MOVEA.L  (A7)+, D0-D2
\end{lstlisting}

The stack itself, the location of {\texttt{A7}} for instance was assigned in the
setup initialization (another tangent to take).

The ARM processor will do this differently.  The author found this simple demonstration:

\begin{lstlisting}
; Prologue - setup
mov     ip, sp                 ; get a copy of sp.
stmdb   sp!, {fp, ip, lr, pc}  ; Save the frame on the stack. See Addendum
sub     fp, ip, #4             ; Set the new frame pointer.
    ...
; Maybe other functions called here.
; Older caller return lr stored in stack frame.
bl      baz
    ...
; Epilogue - return
ldm     sp, {fp, sp, lr}       ; restore stack, frame pointer and old link.
    ...                        ; maybe more stuff here.
bx      lr                     ; return.
\end{lstlisting}

To close off the topic of the Stack, just some more words that help describe the
key features and how they impact development.  Much of this however is handled by
the compiler.  But there are times we don't use the compiler to write a special part
of the code.  Awareness of this is useful for other reasons -- like debugging
stack-overflow, corruption, etc..

Just to get the words, right, I'm taking this from a 
textbook\footnote{{\underline{68000 Assembly Language Programming}}, by 
Gerry Kane, Doug Hawkins, and Lance Leventhal, Osborne/McGraw-Hill, Berkeley CA, 1981, 
page 10-4} the author knows
well, this can apply to just about all fundamental conditioning of the Stack:

\begin{quote}
When passing parameters on the stack, the programmer must implement this
approach as follows:
\begin{enumerate}
\item Decrement the system stack pointer to make room for parameters on the
system stack, and store them using offsets from the stack pointer; or simply
push the parameters on the stack.
\item Access the parameters by means of offsets from the system stack pointer,
remembering that JSR places the return address at the top of the stack.
\item Store the results on the stack by means of offsets from the systems stack
pointer.
\item Clean up the stack before or after returning from the subroutine, so that the
parameters are removed and the results are handled appropriately.
\end{enumerate}
\end{quote}


\item Free-store and Heap.  Here's what Herb Sutter has written about it:

\begin{quote}
Free Store      The free store is one of the two dynamic memory
                 areas, allocated/freed by new/delete.  Object
                 lifetime can be less than the time the storage
                 is allocated; that is, free store objects can
                 have memory allocated without being immediately
                 initialized, and can be destroyed without the
                 memory being immediately deallocated.  During
                 the period when the storage is allocated but
                 outside the object's lifetime, the storage may
                 be accessed and manipulated through a void* but
                 none of the proto-object's nonstatic members or
                 member functions may be accessed, have their
                 addresses taken, or be otherwise manipulated.
\end{quote}

And he also wrote:

\begin{quote}
 Heap            The heap is the other dynamic memory area,
                 allocated/freed by malloc/free and their
                 variants.  Note that while the default global
                 new and delete might be implemented in terms of
                 malloc and free by a particular compiler, the
                 heap is not the same as free store and memory
                 allocated in one area cannot be safely
                 deallocated in the other. Memory allocated from
                 the heap can be used for objects of class type
                 by placement-new construction and explicit
                 destruction.  If so used, the notes about free
                 store object lifetime apply similarly here.
\end{quote}


Now the careful reader of this document will have already objected that this is really
related to C++ and not C!   True.  

But in further research on the very accessible Wikipedia, this very much summarizes the 
points to get across (with apologies for the cut/paste):

\begin{itemize}
\item The C programming language manages memory statically, automatically, or dynamically. 
\item Static-duration variables are allocated in main memory, usually along with the executable code of the program, and persist for the lifetime of the program
\item Automatic-duration variables are allocated on the stack and come and go as functions are called and return.
\item For static-duration and automatic-duration variables, the size of the allocation must be compile-time constant (except for the case of variable-length automatic arrays).
\item Dynamic memory allocation-- memory more explicitly (but more flexibly) managed, typically by allocating it from the free store (informally called the "heap"), an area of memory structured for this purpose.
\end{itemize}
\end{itemize}

So this is where we are left.    In C and for the most part embedded software probably is
developed in C, this point of the Stack, Heap (free-store) is worth discussing with the
candidate and the more they talk about it, the more you will get a sense of their background
facing challenges.

\end{itemize}

\item Comment on this code.  What does it do?  Why?

\begin{lstlisting}
char* ptr;
if (( ptr = (char *) malloc (0)) == NULL )
     puts ( " got a null pointer " );
else
     puts ( " got a valid pointer " );
\end{lstlisting} 

The answer is {\textbf{it depends}}.

\begin{itemize}
\item If the implementation returns \NULL, that's one answer.
\item If the implementation returns a valid pointer, that's another answer.
\end{itemize}

It depends on the implementation of the {\texttt{malloc}} function.

The {\texttt{C}} standard (C17 7.22.3/1)\par
\begin{quote}
If the size of the space requested is zero, the behavior is implementation defined: either a \NULL\ pointer is returned, or the behavior is as if the size were some nonzero value, except that the returned pointer shall not be used to access an object.
\end{quote}

It could return \NULL, or it can return a valid pointer that cannot be dereferenced (but can still be an argument to {\texttt{free}}).

Already we're in a mess. So why do embedded software engineers rely on dynamic memory?  How do embedded software engineers create
the dynamism of allocating resources?  

The candidate should begin to wax poetically about pools and pre-defined buffers, and all sorts of RTOS whistles and bells that may be available,
but then the candidate can also simply say they don't use dynamic memory allocation in real-time embedded software.  Both answers are acceptable,
but we hope that we do hear the candidate talk about both sides of the issue.

About that follow-up question: 
\begin{quote}
Can you explain why I would want to avoid
too many Tasks in an application that uses an RTOS?  How many is too many?
What else should I be concerned about when I have a lot of tasks in an
embedded software application?
\end{quote}

We should look for these in the answer:
\begin{itemize}
\item Tasks that share resources need to be synchronized with an intermediate
mechanisms -- mutex, semaphore, and derivative complex data structures like queues 
or other sorts of ``inter task/process'' communication devices.   
\item Add more tasks, add more code to differentiate between two or more threads of
control vying for the access to a single resource.
\item Add more tasks, add more requirements of each task with space for its
own stack, which puts pressure on the design for a typical embedded software application
that is limited in resources.
\item Add more tasks, add more (a lot more) different ways for faults to occur and the
work and effort to troubleshoot a multi-threaded application (already a challenge
is compounded by more tasks).
\item Fine tuning.  How much quanta given to each task, their priorities, and troubleshooting
starvation of a task from the MCU.
\end{itemize}

The rabbit holes multiply.  The risks multiply and the amount of extra work to ``make
an elegant solution with many threads of control'' devolves into a lot of code just to maintain
the illusion of the elegance.

There's a lot to be said for the single loop, monolithic approach.   Your mileage may vary.

What does your candidate say.  Do the words they say indicate they've been through this
decision process.  Do you want a candidate who hasn't?

\end{enumerate}


\section{Abstraction}

In this section we're going to cover more abstract issues and concepts
that can be asked during the interview for a potential embedded software
engineer (or any software engineer for that matter).

These questions are not specific to embedded software, but there is
a common set of knowledge that software engineers should know before
moving forward in the review process.

I've put into the document a lot more discussion about the answers.
Most of the discussion is my own interpretation of what the answers
should be and also includes a very short primer-like explanation of
why the answer is right.

It is important to know that my answers are not the perfect answers.
\par
Others who have more fine tuned experience with OOP may beg to differ
with the explanation and the answers.  I hope that is just in a few
cases, but I admit that there may be some fundamental differences.
\par
If you feel it is important to correct those differences, let me know.
Your suggestions would be welcome.
\par

Let's begin!

\begin{enumerate}

% top level question
\item What makes a good API?
\par
API is the Application Programming Interface.  Right off the bat, let's
get the terms normalized.  The Application in this case is the software
that uses the API.   For example, the API defined by functions in
{\texttt{stdio.h}} are used by software that expects to use the function
{\texttt{printf}}.   The Interface is the function {\texttt{printf}}
and the Application is some arbitrary code that uses the function.
\par

Another way to view API is how a class is designed.  A class has
methods and data and the methods are what we routinely use to access the data
stored by an instance of the class, or methods are routinely used to
cause behavior implemented by the instance of the class (or by the class
itself in the case of a static method).

\par

In either case, we're interested in the Interface through which actions
and behaviors implemented by some underlying software are exercised.
\par

In light of that basis, the answer to the question ``What makes
a good API?'', it follows that:

\begin{enumerate}
  \item The API should do one thing and do it well.   An example:  {\texttt{stdio.h}} for lack of a better definition is not the best place to define an API
for creating threads, or creating network sockets, etc\...  The {\textit{do 
one thing}} part is that {\texttt{stdio}} does one thing: provide capabilities
for {\textbf{standard}} I/O functionality.   Threads and sockets are not
standard I/O functionality.  They have standards and there are API for
working with threads and sockets, but those functions do not belong in
{\texttt{stdio.h}}.   The ``do it well'' clause means that whatever
the capabilities are, the API provides an  Interface to well designed
and implemented functions for I/O.  A good example is our friend
{\texttt{printf}}.  That function has a lot of features, but few arguments.
It wants a specification for the format of the output ( the first argument ),
and then it wants a variable number of arguments to fit into the format.
It's that simple.  But {\texttt{printf}} does this {\textit{extremely well}}.
\par
When designing an API, the goal should be that the capabilities in the
API are each individually implemented to do their jobs with precision,
and completeness, and as simply as possible.
\item The API should be small as possible, but no smaller.  This goes
to the point earlier about {\textit{do one thing}}.  If the API is kept
focused (as small as possible), then it will naturally converge to
a set of capabilities that are all very close in purpose ({\textit{ do
one thing well}}).  However, the richness of the API should not be
sacrificed for this goal.  Therefore, the clause {\textit{but no smaller}}
means try to keep the API as feature rich as possible while still maintaining
economy of purpose.  If the API is complemented by a function that seems
ever apparently needed along with other functions in the API, then
that function is a good candidate to add to the API.  If a function is
seldom with the other functions in the API -- and moreover -- is likely
to be useful without the API then this is a consideration that the
other function may in fact belong to a different class of functions.
It's up to the designer, but if the designer is sensitive to the
risk of proliferation of functions in the API, these creeping features
are avoided in the target API.
\item Minimize accessibility of Everything.   Not every piece of
data controlled by or used by a class with methods is necessarily
useful to the Programmer.  A simple example is a {\texttt{Coordinate}}
class.  A {\texttt{Coordinate}} might be a class that encapsulates
the information to define a point in some space geometrically.  Each
individual component of data which together defines a Coordinate
may not be of use to a Programmer, but the abstract Coordinate is
useful where the class ( or a peer class ) is capable of performing
functions on Coordinates or functions with Coordinates.  Suppose
the {\texttt{Coordinate}} class was for points in 3-dimensional space.
We would expect each Coordinate to posses an {\texttt{x, y, z}}
member data.  We would not expect the Programmer to want to either
{\texttt{modify}} the individual constituent parts that define
a Coordinate, nor would we even want to or expect the Programmer to
be required to access the constituent parts that define a
Coordinate.  Instead we would probably elect to design the class such
that a given Coordinate is simply itself an argument and the class
provides an API rich enough to make calculations with Coordinates.
Example: Given two Coordinates, what is the Coordinate of the mid-point
between them?  Or, what is the distance between two Coordinates?  Or
given two Coordinates, what is the Coordinate of the destination if
a straight line trajectory for some arbitrary scalar distance is made 
through two Coordinates?  And so on.   All of these types of 
functions of the Coordinate API belong to a Coordinate class probably.
Does the API designer expect to provide access to the internal data
that defines a Coordinate so that the Programmer must implement or
apply yet another API to perform calculations?  Likely, not.  
\item The names of things matter.   When we create the API we are choosing
names for things.  We name functions. We name data held or controlled
by the API too (sometimes).   The class name itself is something we
choose to define.  All of these names are essentially arbitrary and
the Programmer must be aware of them.  They are also names that affect
ways of thinking -- how the functions are used -- when they are used.
For example, consider an API that is responsible for handling special
I/O with a device.  A routine in the API may be a function to indicate that
an operation is completed.   Do we name this function {\texttt{done()}} ?
We could, but what does it mean.  Does it mean that the instance
of the class defining {\texttt{done()}} is no longer useful? Does it mean
that some action (which action?) is completed?  Does it return
truthfulness?   Is the response  to the {\texttt{done()}} function 
true if the work tested by {\texttt{done()}} is completed.  We might
choose names that are simple and abbreviated, but the compiler really
doesn't care how abbreviated the names are, or how compact we try to
fit a name.  What matters is what the Programmer is using -- and the names
of things in the API do matter to the Programmer.   Here are some simple
points to consider when naming things:
   \begin{itemize}
     \item Avoid acronyms.  {\texttt{BRIO}} -- They meant
{\texttt{BlockingRead()}}.  Future users of the API won't know
off hand what {\texttt{BR}} stands for.  Why create confusion? Just write
it out.
     \item For as complex as a function name is, the number of arguments
tends to reduce and vice versa.  If the goal was to have a function for
{\texttt{BlockingRead}} and {\texttt{NonBlockingRead}}, there is a chance
that an argument was avoided -- the flag to indicate if the {\texttt{Read}}
was blocking or not.  So, decide.  What makes more sense?  A tighter
namespace {\texttt{Read()}} but more arguments (flag for blocking or not)?
    \item  Documentation matters\ldots a lot.   It is almost a guarantee
that the author of an API and the author of the implementation of the API
will not be the user of the API.  People change jobs.  People get
transferred. People get bored with whatever they did so they just leave
the code and do something else.  What they leave behind is an API.  If the
API has proper and good documentation, then the next person to pick up the API
can have a very good basis to begin work.  They will understand
what the API is, what the intent is, and they will have a useful basis
to understand how it was supposed to work.  To be clear, in this context
we are referring to the documentation {\textbf{of the API}} not necessarily
the documentation {\textbf{within the implementation of the API}}.  The
best example is already available.  The {\textit{manual page}}.  We use
{\texttt{man}} to look up manual pages for all sorts of API in a system.
Some systems have other file-format schemes for documenting their API.
And still, some systems rely on the documentation to be part of the
API itself -- in the source code headers that define the contour of the API.
In all cases -- manual pages, external documentation like manual pages,
and in-file header/source comments -- these all contribute to the requirement
of describing {\textbf{what}} the API does {\textbf{how}} the API is used.
\par
Many packages exist for assisting with this goal of good API documentation.
In-file markers (think: Doxygen), and other varieties of in-file markings
help automate documentation generation.  Even without the generation
of documentation from in-file markings, the comments in the code itself
(or accompanying MarkDown, or manual pages) are essential for
both the maintainers of the API and the Programmers using the API.
\item Method (function) Design matters\ldots a lot.
   When we look at the API, we find a set of functions that are exposed
   by the API for the Programmer to use.  When those functions (methods)
are designed, we want to maintain order and consistency.  For example
suppose there were two functions in the API that handle processing
data.  One function {\texttt{Result foo(Data d, Size size)}} accepts
a parameter for some data and the size of the data. It returns a
Result.  Another function {\texttt{void bar(Size size, 
Data d, Result\& result)}}
works with data also, it requires the data, the size of the data
and it has an argument that is a reference to store the result. 
\par
The first function returns the Result.  The second function does not
return the Result.  The caller of the first function could theoretically
{\textit{ignore}} the result if there was one.  The second function
demands that the caller provide a variable to store the result.  The
result can still be ignored by the caller, but at least it has to 
have awareness to provide a variable to store it. Further, the second
function switches the order of Size and Data.  The first function, Data
appeared before Size.  In the second function, Data appears after Size.
\par
Despite how pedantic this example is, it just shows something
about design of the API.  Consistency throughout how data is passed into
the functions of the API and consistency throughout how Results are
gained by the caller of the functions of the API.
\par 
API's should be consistent.  If actions taken by the functions in the API
are like-minded (see above), then there is ample reason to choose a set
of Results in a single set so all of the functions in the API return
a value that exists in the set of possible Results.   It doesn't matter
if the nature of Result is a combination of bits that are set for multiple
reasons, it doesn't matter if the Result is a simple native type like
boolean, or some integer value.  At least there is consistency
that any call to any function in the API yields a kind of result in
the returned value.  Also, we look at the formal parameters of the API
methods too.  Are the formal parameters in a consistent form.  For example
Data before Size uniform in the API -- as well as  uniformally 
compared to {\textbf{other API}} in the system.  Is your API an outlier
in terms of names, argument order and return values?   A good API isn't
an outlier in these terms.  A good API ``fits in'' with other API of
the system.
  \end{itemize}
\end{enumerate}  %% What makes a good API question

%% top level question
\item What features in C++ make it object-oriented?

There are several features of object-oriented programming that are
readily available by the language.  Chief among them are
\begin{itemize}
   \item Encapsulation.
   \item Inheritance
   \item Polymorphism.
\end{itemize}

\begin{enumerate}
   \item Encapsulation:   The hiding of data.  Protecting data.
Data (as opposed to methods) that are either hidden away inside the
instance of the class so they are still accessible by methods of the
class, but not exposed outside of the class.  Protection is analogous
to hiding data for a different reason -- two or more classes may
be closely related, each has data, but besides themselves (those friendly
classes), other classes have no rights to the data directly.  And
like protection, the access to the data is metered out by how the
class defines each member data element of the class (just as each
member function is also likewise exposed by access constraints - 
public, private and protected).
  \item Inheritance:  The word speaks for itself.  We use the word
to describe a relationship between two classes.  One class inherits
from another.  What does it inherit?  It inherits two things:  Behavior
and Data.  Specifically, it inherits capabilities (member functions)
at the appropriate level of protection and it also inherits access to
member data (at the appropriate level of protection) in the derived class.
\par
Better terminology is needed.  So for this case we're referring to
the {\textit{base class}} and the {\textit{derived class}}.  In C++
the language provides a way to declare a derived class inheriting from
one or more classes.  We call this {\textit{Multiple Inheritance}}.  Some
other OOP languages don't permit this multiple inheritance.  Those languages
slice the differentiation by declaring a class can {\texttt{extend}} another
class and the class can also {\texttt{implements}} other interfaces.
\par The property of inheritance lets us describe a derived class with the
phrase ``is-a'' relationship.   A {\texttt{Square}} class {\textit{is a}}
{\texttt{Shape}} class.  All squares are shapes.  Not all shapes are 
squares.  Using the phrase ``is-a'' is usually contrasted with another
phrase ``has-a''. A class ``has-a'' member data.   A {\texttt{Shape}}
class has an area.  A {\texttt{Coordinate}} class has a $x, y, z$ value
for the position in 3-dimensional space.
\par
But, back to the point of Inheritance.  If the goal is simply that the
derived class mimics exactly the same behavior of the base class(es). Then
inheritance is not useful.  What is useful is to look at a class
in terms of the API it provides when thinking about inheritance.  The 
methods assigned to a base class form a basis of functional behavior.
If the derived class inherits from the base class, the derived class is
free to override those methods.  The derived class can provide different
implementation to the methods that are inherited from the base class.
For example, the base class {\texttt{Shape}} may have a method called
{\texttt{float area()}}.  The intent is that this method will return
the calculated area of the shape.   We have derived two classes from
{\texttt{Shape}} -- {\texttt{Square}} and {\texttt{Circle}}.  If
we wanted to compute the area of the square object or compute the area of
the circle object, the algorithm is different.  For a square with side
length $s$, the area is $s^2$.   A circle's area is $\pi r^2$ for circle
radius $r$.  We know this, but the {\texttt{Shape}} base class does not
and cannot know this.  Therefore we design the classes so that
a {\texttt{Shape}} class has a method {\texttt{float area())}} but does
not implement it (it cannot know how).  But any derived class of
{\texttt{Shape}} must implement it.   So in some sense, the derived
classes inherit the {\textit{interface}} declared in {\texttt{Shape}}.
\item Polymorphism:  This is related to Inheritance. Without inheritance,
polymorphism is really not possible.  The word means that the instance
of an object takes on the role of a different class depending on which
interface to the object is available.
The definition from Bjarne Stroustrup is as follows:
\begin{quote}
polymorphism - providing a single interface to entities of different types. Virtual functions provide dynamic (run-time) polymorphism through an interface provided by a base class. Overloaded functions and templates provide static (compile-time) polymorphism. TC++PL 12.2.6, 13.6.1, D\&E 2.9.
\end{quote}
\par
For example, back to our simple shapes scenario.   We might have
declared that the {\texttt{Shape}} class is an abstract class. This means
that at least one of the methods in the class is declared purely virtual.
Abstract classes are by their definition meant to be a base class.
Our method {\texttt{float area()}} in {\texttt{Shape}} is specified
so that any derived class provides the method to compute the area of the
derived class.   The method is declared purely virtual:

\begin{lstlisting}
class Shape
{
   public:
      virtual float area() const = 0;

   // ...
};
\end{lstlisting}

A class derived from {\texttt{Shape}} looks like this:


% Listing of the Shape class
\begin{lstlisting}
\end{lstlisting}

And finally, a simple square class:

% Listing of the Square class
\begin{lstlisting}
\end{lstlisting}

How polymorphism can be used.

Suppose we make some objects:

\begin{lstlisting}
Ellipse* pEllipse = new Ellipse(1.0f, 1.5f);
\end{lstlisting}

This is a pointer to a new object of type {\texttt{Ellipse}}. We
constructed it so that the minor and major axis lengths are
$1.0, 1.5$ respectfully.

Another class might accept inputs of various shapes.  This function is
not aware of all the kinds of shapes that exist. But, it only has to 
verify that the area of the shape passed to the function is within a certain
limit:

\begin{lstlisting}
class ShapeTester
{
    public:
       ShapeTester(float space)
       {
            maxSpace = space
       }

       // If the shape fits the allowed space
       //  return true
       // Otherwise return false;
   
       bool shapeFits(Shape* pShape)
       {
            float area = pShape->area();
            bool fits = (area <= maxSpace);
            return fits;
       }

    private:
       float maxSpace;
};
\end{lstlisting}

Right away, we see in the API of the {\texttt{ShapeTester}} that
it wants a pointer to a {\texttt{Shape}}, not a pointer to
some arbitrary shape like {\texttt{Square}}, {\texttt{Circle}},
or {\texttt{Ellipse}}.

In other words -- the tester is ambivalent about the shape. All it
wants is the Interface to an object to get the area of the shape.

We made all of our shape classes inherit {\texttt{Shape}} so
no matter which derived class we made from {\texttt{Shape}} we are
sure that the tester will be able to get the area of the shape.

Back to our example code:

\begin{lstlisting}
   ShapeTester* pShapeTester = new ShapeTester(10.0f);
   // ...

   Ellipse* pEllipse = new Ellipse(1.0f, 1.5f);
   // ...

   bool fits = pShapeTester->shape_fits(pEllipse);

   if (fits) {
        std::cout << "Shape fits" << std::endl;
   }
   else {
        std::cout << "Shape does not fit" << std::endl;
   }
\end{lstlisting}


What happens here? Polymorphism.  The class {\texttt{Ellipse}} can be
viewed as an object of type {\texttt{Shape}} and it can be viewed
as an object of type {\texttt{Ellipse}}.   One instance of 
{\texttt{Ellipse}} exists, but there are two interfaces to the object:
The {\texttt{Shape}} interface and the {\texttt{Ellipse}} interface.

\par
The {\texttt{Shape}} interface provides access to all of the methods
and data accessible that are in the base class {\texttt{Shape}}.
The {\texttt{Ellipse}} interface provides access to all of the
methods and data accessible that are in the derived class {\texttt{Ellipse}}.
\par 
What happens is equivalent as if the code was this:

\begin{lstlisting}
   // ...

   Shape* pShapeInterface = (Shape*) pEllipse;

   bool fits = pShapeTester->shape_fits(pShapeInterface);
   // ...
\end{lstlisting}

The point is that the object {\texttt{pEllipse}} two interfaces.
One of the interfaces is of type {\texttt{Shape}} and the other interface
is of type {\texttt{Ellipse}}.  It's actually a pointer {\textbf{to}}
a {\texttt{Shape}}, hence {\texttt{Shape*}}.
\par
It would be complicated and overly rigid if the {\texttt{ShapeTester}}
had to be written with overloaded functions, one per shape.  In line
with what Stroustrup wrote, the single method in the {\texttt{ShapeTester}}
will work with {\textbf{any}} derived class from {\texttt{Shape}}.
\par

This whole example is a very simple way of looking at Polymorphism.
It relies on the ``is-a'' property of a class hierarchy.  It relies on
the {\texttt{Ellipse}} class is derived {\textbf{from}} a {\texttt{Shape}}.
\par
But, there are far more elaborate ways to use the concept of Polymorphism.

Here's a short list:

\begin{itemize}
\item Operator overloading.  Suppose a set of classes represents numerical
quantities, but the default operators provided aren't aware of how to
perform calculations on these new classes.   Adding a numerical quantity of
type {\texttt{Foo}} to another numerical quantity of type {\texttt{Bar}}.
(assuming that {\texttt{Foo}}, and {\texttt{Bar}} are just classes
with backing data that is numerical) would be difficult unless we can
overload the basic operators $+, -, *, /$, etc\ldots.
\item A more specific Polymorphism example has to do with abstractions
and patterns like {\textbf{class factories}}, or even {\textbf{COM}}
models.  This is where the idea that a derived class ``is-a'' becomes
more subtle.

In the case of an Abstract Class Factory, we choose to design a class
capable of returning an instance (an object) that is a certain class
based solely on some argument applied to the Factory.    It's analogous
to asking for a burger, or asking for fries, or a milkshake -- either one
but asking the same person for it.  The person is the Class Factory,
and the item requested is the instance of something class {\texttt{Burger}},
{\texttt{Fries}}, or {\texttt{MilkShake}}.  The Factory has to
have prior knowledge of the classes it is capable of creating object
instances.
\par
The difference between this and COM (among many) is that the derived
classes using COM are written to define at compile-time the hierarchical
dependencies of a derived class.  But {\textit{accessing the interface}}
to a specific base class of a derived class depends on the request knowing
some information (using a UUID) which provides the {\texttt{IUnknown}}
interface method {\texttt{QueryInterface}} information to cast
a result parameter based on the UUID.
\end{itemize}

In other words, the coupling of a derived class to its base class can
be very tight (inheritance based alone).  Or, it can be loosely
coupled like COM or other mechanisms where the requester merely just
asks for an interface and the grantor creates a pointer to an interface
of some other class.

  \end{enumerate}


\item How is data hidden from other classes in C++ ?

Patterns: Passkey, PIMPL\\
Language: {\texttt{private, protected, public}} \\
\par


As far as I'm aware there are language features and software features
to do this.

In the language, within a class defined in the language, members can be
restricted or unrestricted from other classes by assigning
 {\texttt{private, protected, public}} to that member.

A class can access all of its members.  And:

\begin{itemize}
\item {\texttt{public}} - Any other class can access the public member.
\item {\texttt{protected}} - Only derived classes can access the protected member. Or classes that are {\texttt{friend}} to the class.
\item {\texttt{private}} - No other class can access the private member,
except those classes that are a {\texttt{friend}} to the class.
\end{itemize}


Example:
\begin{lstlisting}
class Alpha
{
     public:
          int x;
          void func();
}

// ...

Alpha a;

\end{lstlisting}

Any other class may access {\texttt{a.x}}.  Any other class
may call the function {\texttt{a.func()}}.  




%% Section
\end{enumerate} 


Other quotes to  use:

Note: Old line C programmers sometimes don't like references since they provide reference semantics that isn't explicit in the caller's code. After some C++ experience, however, one quickly realizes this is a form of information hiding, which is an asset rather than a liability. E.g., programmers should write code in the language of the problem rather than the language of the machine.


John Carmack:

NULL pointers are the biggest problem in C/C++, at least in our code. The dual use of a single value as both a flag and an address causes an incredible number of fatal issues. C++ references should be favored over pointers whenever possible; while a reference is “really” just a pointer, it has the implicit contract of being not-NULL. Perform NULL checks when pointers are turned into references, then you can ignore the issue thereafter.



\end{document}

